\chapter{BUSINESS RULES VÀ KPI REPORTING RULES}

\noindent\textbf{Mục tiêu chương}

\begin{enumerate}[label=-]
\item Tóm tắt các business rules cốt lõi đang chi phối hệ thống báo cáo năng lượng.
\item Ghi nhận rõ source-of-truth rules, period rules, display rules, và KPI coverage rules.
\item Giúp kỹ sư phát triển, người rà soát nghiệp vụ, và người vận hành có cùng cách hiểu về logic hiện hành.
\item Tạo một chapter chuẩn để mở rộng thành tài liệu rules chính thức khi dự án tiếp tục hoàn thiện.
\end{enumerate}

\section{Vai trò của business rules trong dự án}

Đối với hệ thống báo cáo năng lượng, business rules không chỉ quyết định cách hiển thị báo cáo mà còn quyết định cách dữ liệu được chọn, được tổng hợp, và được chấp nhận để phát hành. Các rule được mô tả trong chapter này phải luôn bám sát implementation thực tế và nhất quán với các tài liệu nguồn như \texttt{docs/business\_rules.md} và \texttt{docs/kpi\_reporting\_rules.md}.

Trong các hệ thống ETL hoặc reporting pipeline, việc duy trì đúng grain của dữ liệu, đúng source-of-truth, và đúng điều kiện tổng hợp là điều kiện tiên quyết để tránh sai lệch kết quả \cite{ref_kimball}. Với dự án này, phần rules là lớp bảo vệ quan trọng cho tính đúng đắn của output cuối cùng.

\section{General reporting invariants}

Các quy tắc tổng quát sau áp dụng cho toàn bộ report stack:

\begin{enumerate}[label=-]
\item báo cáo được tổ chức thành các section có cấu trúc rõ ràng, tối thiểu gồm Electricity, Utility, KPI, và phần metadata liên quan;
\item mỗi report phải hỗ trợ so sánh giữa kỳ hiện tại và kỳ trước tương ứng;
\item runtime hiện tại phát hành các kỳ \textit{daily}, \textit{weekly}, và \textit{monthly}; rules nền vẫn phải tương thích với các custom range trong tương lai nếu chúng được bật lại;
\item các bảng daily phải luôn dense, không bỏ mất ngày chỉ vì thiếu dữ liệu;
\item giá trị bằng 0 là hợp lệ, còn giá trị thiếu phải hiển thị rõ bằng dấu \texttt{-};
\item logic nghiệp vụ phải nằm ở backend/service layer, không nhân bản trong template HTML hay PDF.
\end{enumerate}

Đối với mọi metric có so sánh, backend cần chuẩn bị đầy đủ current value, previous value, delta, và delta percentage theo công thức sau:

\begin{equation}
	\Delta = \text{Current Value} - \text{Previous Value}
	\label{eq:delta_value}
\end{equation}

\begin{equation}
	\Delta\% = \begin{cases}
		\frac{\Delta}{\text{Previous Value}} \times 100, & \text{khi Previous Value} \neq 0 \\
		\texttt{-}, & \text{khi Previous Value} = 0
	\end{cases}
	\label{eq:delta_pct}
\end{equation}

Cách xử lý này giúp report phản ánh trung thực sự thay đổi mà không tạo ra phần trăm ảo trong trường hợp mẫu số bằng 0.

\section{Source of truth và ownership map}

Bảng \ref{tab:rule_ownership_matrix} tóm tắt nguồn dữ liệu chính thức và lớp chịu trách nhiệm diễn giải chúng.

\begin{table}[h]
	\centering
	\renewcommand{\arraystretch}{1.3}
	\caption{Rule ownership and source-of-truth matrix}
	\label{tab:rule_ownership_matrix}
	\begin{tabularx}{\linewidth}{|>{\RaggedRight\arraybackslash}p{3.1cm}|>{\RaggedRight\arraybackslash}p{3.2cm}|>{\RaggedRight\arraybackslash}X|}
		\hline
		\textbf{Domain} & \textbf{Primary Source / Owner} & \textbf{Critical Rule} \\ \hline
		Electricity official totals & \texttt{energy\_kpi} & Plant totals và area official totals phải dùng giá trị đã được phê duyệt, không tự cộng lại từ raw energy views. \\ \hline
		KPI values and production-linked reporting & \texttt{energy\_kpi} & KPI reporting chỉ dùng các block đã tính sẵn và đã được lưu, tuân thủ coverage-first logic. \\ \hline
		Utility daily usage & utility usage source & Utility summary và daily detail phải so sánh current/previous và giữ dense rows cho khoảng ngày cần hiển thị. \\ \hline
		Sensor monitoring & \texttt{processvalue} + sensor metadata mapping & Raw sensor rows phải được aggregate thành payload có chất lượng dữ liệu rõ ràng, hỗ trợ anomaly-ready analysis. \\ \hline
		Presentation rules & backend services + render context & Template chỉ nhận dữ liệu đã chuẩn hóa từ \texttt{report\_context}; không tái tính nghiệp vụ trong view/PDF. \\ \hline
	\end{tabularx}
\end{table}

Source-of-truth discipline này rất quan trọng vì nếu sai nguồn dữ liệu hoặc sai grain tổng hợp, toàn bộ report có thể trở nên nhất quán về hình thức nhưng sai về bản chất. Đây cũng là nguyên tắc nền trong các hệ thống dữ liệu định hướng báo cáo và warehouse \cite{ref_kimball,ref_mysql}.

\section{Electricity rules}

\subsection{Official totals và residual load}

Các rule điện năng cốt lõi gồm:

\begin{enumerate}[label=-]
\item official plant total và official area total phải lấy từ \texttt{energy\_kpi};
\item không được suy ra official totals bằng cách cộng toàn bộ meter values từ energy views, vì các energy views có thể chứa main feeders và downstream overlap;
\item main feeder definitions và area topology phải được quản lý qua metadata cấu hình của dự án;
\item residual load được xem như một virtual meter khi có đủ điều kiện tính hợp lệ.
\end{enumerate}

Residual load được định nghĩa như sau:

\begin{equation}
	\text{Residual Load} = \text{Official Area Total} - \sum \text{Branch Meters (excluding main feeders)}
	\label{eq:residual_load}
\end{equation}

Residual load đại diện cho phần tải chưa được đo trực tiếp bởi branch submeters nhưng vẫn thuộc về area official total, do đó cần được xử lý nhất quán như một phần hợp lệ của downstream reporting logic.

\subsection{Ranking, summary, và daily detail}

Các quy tắc hiển thị Electricity hiện hành gồm:

\begin{enumerate}[label=-]
\item top meter ranking được xếp theo current period consumption;
\item ranking phải đính kèm previous-period value để hỗ trợ so sánh;
\item main distribution feeders phải bị loại khỏi Top Meter ranking và Top 1 daily meter logic;
\item residual load có thể xuất hiện như một virtual meter nếu giá trị hợp lệ;
\item display limit hiện hành là Top 10 cho daily và weekly, Top 5 cho monthly;
\item daily summary phải bao gồm total energy, top 1 meter, active meter count, và average per active meter;
\item daily detail phải render full list of configured meter columns, kể cả khi có ngày hoặc cột không có dữ liệu.
\end{enumerate}

Ngoài ra, daily detail của Electricity phải cho phép backend chuẩn bị sẵn heatmap-related metadata để template chỉ cần render, thay vì tự quyết định logic tô màu hoặc so sánh tại tầng view.

\section{Utility và sensor monitoring rules}

\subsection{Utility reporting rules}

Utility domain bao phủ các nhóm tiêu thụ như nước, chilled water, compressed air, steam, và các utility metrics khác đã được cấu hình. Các rule chính gồm:

\begin{enumerate}[label=-]
\item utility summary phải hỗ trợ current, previous, delta, và delta percentage;
\item daily utility detail phải render full date range, không bỏ sót ngày;
\item mỗi cột utility phải dùng business-level naming nhất quán thay vì lộ thẳng raw sensor naming;
\item missing values phải hiển thị bằng \texttt{-}, không để trống hoặc gây hiểu nhầm là 0.
\end{enumerate}

\subsection{Sensor monitoring rules}

Sensor monitoring được xây dựng từ raw sensor data trong \texttt{processvalue}. Mỗi metric cần có mapping rõ giữa business key, display name, unit, và source sensor column. Sau khi aggregate, payload nên bảo toàn các tín hiệu chất lượng dữ liệu sau:

\begin{enumerate}[label=-]
\item min value;
\item avg value;
\item max value;
\item latest value;
\item sample count;
\item non-null count;
\item zero count;
\item negative count;
\item anomaly-ready quality signals như coverage, zero-heavy patterns, và negative-tolerance checks.
\end{enumerate}

Bộ output hiện tại của sensor monitoring phải đủ linh hoạt để phục vụ cả daily và periodic Utility surfaces, bao gồm các family như \texttt{metric\_columns}, \texttt{daily\_rows}, \texttt{overview\_cards}, \texttt{groups}, \texttt{anomaly\_rows}, và \texttt{period\_detail\_tables}.

\section{KPI reporting rules}

\subsection{KPI definition và source-of-truth rule}

KPI của dự án đang được dùng theo đơn vị \texttt{kWh/Ton}. Giá trị KPI được dẫn xuất từ năng lượng và sản lượng, nhưng điều quan trọng là report không được tự recompute KPI từ raw electricity views. KPI reporting phải sử dụng các row đã được tính sẵn và lưu trong \texttt{energy\_kpi}.

Các nguyên tắc bất biến gồm:

\begin{enumerate}[label=-]
\item coverage-first logic là bắt buộc;
\item không prorate Week, Month, hoặc Year rows thành các giá trị daily nhân tạo;
\item không split broader blocks thành synthetic rows;
\item nếu có nhiều row cho cùng một reporting block, luôn chọn row có \texttt{dt\_lastupdate} mới nhất;
\item KPI rendering phải phản ánh stored business-approved values, không phản ánh giá trị suy diễn.
\end{enumerate}

\subsection{Coverage hierarchy và selection rules}

Coverage hierarchy hiện hành là:

\begin{center}
	\textbf{Day $>$ Week $>$ Month $>$ Year}
\end{center}

Nghĩa là hệ thống luôn ưu tiên block chi tiết hơn khi block đó hợp lệ. Một broader block chỉ được dùng như chính block đã lưu của nó, không được chia nhỏ lại để phủ lên các ngày thiếu.

Bảng \ref{tab:kpi_block_rules} tóm tắt điều kiện chấp nhận hoặc loại bỏ các KPI blocks.

\begin{table}[h]
	\centering
	\renewcommand{\arraystretch}{1.3}
	\caption{KPI block selection rules}
	\label{tab:kpi_block_rules}
	\begin{tabularx}{\linewidth}{|>{\RaggedRight\arraybackslash}p{2.2cm}|>{\RaggedRight\arraybackslash}X|>{\RaggedRight\arraybackslash}X|}
		\hline
		\textbf{Block Type} & \textbf{Accepted When} & \textbf{Reject When} \\ \hline
		Day & row khớp chính xác một ngày cần hiển thị & row trùng nhưng có \texttt{dt\_lastupdate} cũ hơn block mới hơn \\ \hline
		Week & block tuần nằm trọn trong uncovered range và không chồng lên các ngày đã có Day coverage hợp lệ & block chỉ overlap một phần hoặc cần bị chia nhỏ để dùng \\ \hline
		Month & block tháng được dùng như full stored block hợp lệ cho kỳ tháng tương ứng & block tháng chỉ fit một phần custom range hoặc cần prorating \\ \hline
		Year & block năm chỉ dùng như full stored block khi thật sự là kỳ hợp lệ tương ứng & block năm làm override coverage chi tiết hơn hoặc cần phân rã thành smaller units \\ \hline
	\end{tabularx}
\end{table}

\subsection{Coverage result và daily detail}

Mỗi KPI result cần expose tối thiểu các trường:

\begin{enumerate}[label=-]
\item \texttt{coverage\_display};
\item \texttt{coverage\_note};
\item \texttt{uncovered\_ranges};
\item \texttt{is\_complete}.
\end{enumerate}

Daily KPI detail phải là một dense validation surface. Ngay cả khi ngày đó không có KPI coverage hợp lệ, hàng vẫn phải tồn tại và hiển thị trạng thái \textit{Missing} hoặc \textit{Partial} cùng các dấu \texttt{-} cho production, energy, và KPI values chưa có.

Triết lý hiển thị ở đây là: \textbf{partial nhưng đúng vẫn tốt hơn complete nhưng sai}. Điều này cũng phù hợp với tinh thần quản trị dữ liệu và quản lý hiệu suất năng lượng có kiểm soát \cite{ref_iso50001}.

\section{Data quality và implementation guardrails}

Ngoài các rule nghiệp vụ thuần túy, hệ thống còn cần các guardrail kỹ thuật sau:

\begin{enumerate}[label=-]
\item không crash vì bad records hoặc null-heavy rows;
\item các bước fetch, aggregate, render phải có logging đủ để debug;
\item nên log row counts và key execution facts cho các bước quan trọng;
\item tránh hard-coded logic khi có thể biểu diễn bằng metadata hoặc config;
\item giữ code ở dạng modular, testable, và dễ rà soát;
\item khi docs và implementation mâu thuẫn, phải báo mismatch thay vì tự suy diễn một phía là đúng.
\end{enumerate}

\section{Checklist rà soát rules trước khi phát hành}

Trước khi phát hành report hoặc technical PDF có liên quan đến logic nghiệp vụ, nên rà soát tối thiểu các câu hỏi sau:

\begin{enumerate}[label=-]
\item source dữ liệu cho từng khối đã đúng source-of-truth chưa;
\item KPI blocks có tuân thủ coverage-first và no-prorating rules chưa;
\item delta và delta percentage có được tính an toàn khi previous value bằng 0 không;
\item dense row behavior có được giữ đúng ở daily detail và KPI detail không;
\item template có đang chỉ render dữ liệu hay đã lỡ mang business logic xuống tầng presentation;
\item các thay đổi về metadata, feeder exclusions, hay sensor mappings đã được phản ánh vào docs và report output chưa.
\end{enumerate}

\section{Tổng kết chương}

Chapter này xác định rõ lớp rules cốt lõi của dự án: source-of-truth discipline, electricity rules, utility/sensor monitoring rules, KPI coverage rules, và các guardrail kỹ thuật để hệ thống luôn cho ra report đúng và có thể bảo trì. Ở trạng thái hiện tại, bộ tài liệu cũng đã có chapter riêng cho output structure và release workflow, deployment and operations, troubleshooting, và release-readiness checklist.
